\chapter*{CHƯƠNG 3: PHÂN TÍCH VÀ THIẾT KẾ HỆ THỐNG}
\phantomsection
\addcontentsline{toc}{chapter}{CHƯƠNG 3: PHÂN TÍCH VÀ THIẾT KẾ HỆ THỐNG}

Chương này trình bày thiết kế hệ thống \projectname{} ở mức kiến trúc, dữ liệu, API và luồng xử lý. Mục tiêu của chương là cho thấy hệ thống được thiết kế như một phần mềm có cấu trúc, không phải chỉ là một trang web upload file đơn giản.

\section*{3.1. Kiến trúc tổng quan}
\phantomsection
\addcontentsline{toc}{section}{3.1. Kiến trúc tổng quan}

Hệ thống được chia thành ba vùng chính: trình duyệt người dùng, ứng dụng Next.js và tầng lưu trữ. Trình duyệt chịu trách nhiệm tương tác ví và mã hóa/giải mã. Backend chịu trách nhiệm xác thực, phân quyền, quản lý metadata và audit. Storage chỉ lưu ciphertext.

\begin{figure}[H]
\centering
\includegraphics[width=\textwidth]{architecture.png}
\caption{Kiến trúc tổng quan của hệ thống}
\label{fig:architecture}
\end{figure}

Cách phân tách này có các ưu điểm:

\begin{itemize}
  \item Logic nhạy cảm về nội dung file nằm ở client.
  \item Backend tập trung vào quyền truy cập và vòng đời tài liệu.
  \item Storage có thể thay thế giữa local filesystem và Cloudflare R2.
  \item Dễ kiểm thử từng tầng: UI, API, database và storage.
\end{itemize}

\section*{3.2. Thiết kế module}
\phantomsection
\addcontentsline{toc}{section}{3.2. Thiết kế module}

\begin{table}[H]
\centering
\caption{Các module chính trong hệ thống}
\begin{tabularx}{\textwidth}{|p{3.5cm}|X|}
\hline
\textbf{Module} & \textbf{Trách nhiệm} \\ \hline
Auth module & Tạo nonce, xác minh chữ ký ví, tạo JWT session và đăng xuất. \\ \hline
Encryption identity module & Lưu public key của ví để phục vụ key envelope và threshold share. \\ \hline
File module & Lưu metadata file, version, access mode, vault và trạng thái hủy. \\ \hline
Storage module & Upload, lấy, kiểm tra và xóa ciphertext từ local hoặc R2. \\ \hline
Token module & Cấp, kiểm tra, liệt kê và thu hồi token truy cập. \\ \hline
Vault module & Quản lý kho nhóm, thành viên và vai trò. \\ \hline
Approval module & Quản lý yêu cầu phê duyệt, contribution và approval token. \\ \hline
Audit module & Ghi và truy vấn các sự kiện quan trọng. \\ \hline
UI module & Dashboard, upload wizard, downloader, viewer và thông báo. \\ \hline
\end{tabularx}
\end{table}

\section*{3.3. Thiết kế cơ sở dữ liệu}
\phantomsection
\addcontentsline{toc}{section}{3.3. Thiết kế cơ sở dữ liệu}

Hệ thống sử dụng SQLite cho metadata. Dữ liệu file thực tế không lưu trong database mà lưu ở storage dưới dạng ciphertext, còn database lưu cid và thông tin quyền.

\small
\begin{longtable}{|p{4.6cm}|p{8.5cm}|}
\caption{Các bảng dữ liệu chính} \\
\hline
\textbf{Bảng} & \textbf{Vai trò} \\ \hline
\seqsplit{files} & Lưu metadata tài liệu: owner, cid, tên file, mime, kích thước, iv, access mode, version, vault, trạng thái hủy. \\ \hline
\seqsplit{tokens} & Lưu token truy cập, file tương ứng, người nhận, thời hạn, trạng thái thu hồi và mục đích token. \\ \hline
\seqsplit{encryption\_identities} & Lưu public key của ví và chữ ký xác nhận public key thuộc về ví đó. \\ \hline
\seqsplit{key\_envelopes} & Lưu key envelope cho token đích danh. \\ \hline
\seqsplit{vaults} & Lưu thông tin kho nhóm. \\ \hline
\seqsplit{vault\_members} & Lưu thành viên kho và vai trò owner/editor/viewer. \\ \hline
\seqsplit{vault\_threshold\_policies} & Lưu chính sách phê duyệt ngưỡng của kho. \\ \hline
\seqsplit{threshold\_file\_shares} & Lưu các share đã mã hóa cho từng thành viên. \\ \hline
\seqsplit{approval\_requests} & Lưu yêu cầu truy cập của người nhận. \\ \hline
\seqsplit{approval\_contributions} & Lưu đóng góp phê duyệt của từng người duyệt. \\ \hline
\seqsplit{audit\_events} & Lưu nhật ký hành động quan trọng; có trigger ngăn sửa/xóa. \\ \hline
\seqsplit{rate\_limits} & Lưu bộ đếm giới hạn request cho API nhạy cảm. \\ \hline
\end{longtable}
\normalsize

Một số quan hệ quan trọng:

\begin{itemize}
  \item Một file có thể có nhiều token.
  \item Một vault có nhiều thành viên và một chính sách phê duyệt.
  \item Một yêu cầu phê duyệt thuộc về một file và một requester.
  \item Một yêu cầu có nhiều contribution từ các người duyệt khác nhau.
  \item Audit event gắn với actor, action, resource type và resource id.
\end{itemize}

\section*{3.4. Thiết kế API}
\phantomsection
\addcontentsline{toc}{section}{3.4. Thiết kế API}

\small
\begin{longtable}{|p{1.8cm}|p{4.5cm}|p{6.8cm}|}
\caption{Nhóm API chính} \\
\hline
\textbf{Phương thức} & \textbf{Endpoint} & \textbf{Mục đích} \\ \hline
POST & \seqsplit{/api/auth/start} & Sinh thông điệp nonce để ví ký. \\ \hline
POST & \seqsplit{/api/auth/verify} & Xác minh chữ ký ví và tạo session. \\ \hline
POST & \seqsplit{/api/storage/upload} & Nhận ciphertext và lưu vào provider storage. \\ \hline
GET & \seqsplit{/api/storage/get} & Trả ciphertext nếu token hoặc approval hợp lệ. \\ \hline
POST & \seqsplit{/api/files} & Tạo metadata file, token và share nếu có. \\ \hline
DELETE & \seqsplit{/api/files/[id]} & Chủ sở hữu hủy tài liệu. \\ \hline
POST & \seqsplit{/api/tokens/issue} & Cấp token mới cho file. \\ \hline
POST & \seqsplit{/api/tokens/validate} & Kiểm tra token trước khi mở file. \\ \hline
POST & \seqsplit{/api/tokens/revoke} & Thu hồi token. \\ \hline
GET & \seqsplit{/api/approvals} & Liệt kê yêu cầu phê duyệt liên quan đến người dùng. \\ \hline
POST & \seqsplit{/api/approvals/[id]/approve} & Gửi contribution phê duyệt cho yêu cầu. \\ \hline
GET & \seqsplit{/api/audit} & Truy vấn audit event được phép xem. \\ \hline
GET & \seqsplit{/api/health} & Kiểm tra database và provider storage. \\ \hline
\end{longtable}
\normalsize

\section*{3.5. Thiết kế luồng upload và chia sẻ}
\phantomsection
\addcontentsline{toc}{section}{3.5. Thiết kế luồng upload và chia sẻ}

Luồng upload gồm các bước:

\begin{enumerate}
  \item Người dùng kết nối ví và có session hợp lệ.
  \item Người dùng chọn file trong upload wizard.
  \item Trình duyệt tạo khóa mã hóa và mã hóa file bằng Web Crypto.
  \item Client upload ciphertext lên \code{/api/storage/upload}.
  \item Storage trả về \code{cid}.
  \item Client gửi metadata, iv, access mode và key material hợp lệ tới \code{/api/files}.
  \item Backend lưu metadata, tạo token ban đầu và ghi audit event.
\end{enumerate}

Điểm quan trọng là file rõ chỉ tồn tại tại trình duyệt. Backend không nhận file gốc trong thiết kế chủ đích.

\section*{3.6. Thiết kế luồng phê duyệt A+B cho C}
\phantomsection
\addcontentsline{toc}{section}{3.6. Thiết kế luồng phê duyệt A+B cho C}

Luồng phê duyệt được thiết kế để minh họa kiểm soát truy cập nhiều người. Hình \ref{fig:approvalseq} mô tả trình tự chính.

\begin{figure}[H]
\centering
\includegraphics[width=\textwidth]{approval-sequence.png}
\caption{Sequence diagram luồng phê duyệt A+B cho C}
\label{fig:approvalseq}
\end{figure}

Các quy tắc chính:

\begin{itemize}
  \item Người yêu cầu không được tự phê duyệt yêu cầu của mình.
  \item Contribution của mỗi người duyệt chỉ được tính một lần.
  \item Token approval chỉ được tạo khi số contribution hợp lệ đạt ngưỡng.
  \item Token approval gắn với requester, không phải token mở tự do cho mọi ví.
\end{itemize}

\section*{3.7. Thiết kế giao diện}
\phantomsection
\addcontentsline{toc}{section}{3.7. Thiết kế giao diện}

Giao diện gồm bốn màn hình chính: giới thiệu, upload, dashboard và viewer. Các màn hình dùng chung hệ nhận diện, chuyển động và logo của dự án.

\clearpage
\begin{figure}[H]
\centering
\includegraphics[width=\textwidth,height=0.74\textheight,keepaspectratio]{live-home.png}
\caption{Trang giới thiệu hệ thống}
\end{figure}

\begin{figure}[H]
\centering
\includegraphics[width=\textwidth,height=0.74\textheight,keepaspectratio]{live-upload.png}
\caption{Upload wizard và thiết lập quyền chia sẻ}
\end{figure}

\begin{figure}[H]
\centering
\includegraphics[width=\textwidth,height=0.74\textheight,keepaspectratio]{live-dashboard.png}
\caption{Dashboard quản lý tài liệu, kho và yêu cầu}
\end{figure}

\begin{figure}[H]
\centering
\includegraphics[width=\textwidth,height=0.74\textheight,keepaspectratio]{live-download.png}
\caption{Luồng nhận tài liệu và protected viewer}
\end{figure}

\clearpage
